\chapter{Incident response and validation}
\label{ch:ch4}

The current response prototype receives Falco syscall events through Falcosidekick
over HTTP. Falcosidekick supplies a shared token to the handler. Separate ingress
policies restrict Falco-to-Falcosidekick and Falcosidekick-to-handler traffic. These
controls assume a trusted cluster administration boundary and do not replace TLS,
durable event delivery or a production controller.

\section{Current deployment procedure}

Use the repository README from the repository root. It pins Falco chart 9.1.0,
uses modern eBPF (requiring kernel BTF), creates the authentication Secret and applies
the Cilium isolation policy before enabling automatic response. Chapter~\ref{ch:ch2}
describes the same deployment sequence and its environment requirements.

The handler Service listens on port 80 and forwards to container port 5000.
The Deployment uses a non-root user, read-only root filesystem, dropped capabilities,
RuntimeDefault seccomp, resource limits and health/readiness endpoints. Its Role is
limited to pod operations in the default namespace. Workload label checks are
application checks: Kubernetes RBAC itself is not label-scoped.

When migrating, inspect and remove obsolete cluster-scoped bindings only after
confirming ownership. Applying a namespaced Role does not remove a ClusterRoleBinding.
The original names were \texttt{falco-handler-deployment-binding} and
\texttt{falco-deployment-manager}; the intermediate revision used
\texttt{falco-handler-pod-binding} and \texttt{falco-pod-manager}.

\section{Response sequence}

The handler validates the token, JSON, syscall source, exact rule name and namespace
before reading the target pod. It then validates the workload label. The accepted
rules are the official reverse-shell rule ending in ``in Container'', PTRACE anti-debug
attempt, and the custom network-tools rule.

The handler removes only the configured selector labels (app and pod-template-hash),
preserving unrelated metadata. It adds the quarantine and manager labels and a UTC
quarantine timestamp. The observed resourceVersion detects stale updates.

\begin{lstlisting}[language=bash]
# Inspect the result; this does not itself prove network enforcement.
kubectl -n default get pods --show-labels
kubectl -n default get endpointslices \
  -l kubernetes.io/service-name=vulnapp
\end{lstlisting}

The Service and ReplicaSet reconcile asynchronously: the pod leaves the service
selector and a replacement can be created. Two replicas reduce disruption during
a single incident, but no zero-downtime guarantee is made. The PDB applies to voluntary
eviction, not to response-driven label changes.

\section{Network quarantine}

The CiliumNetworkPolicy selects the quarantine label and uses explicit ingressDeny
and egressDeny entries covering all entities. Empty ingress/egress lists must not be
used as an assumed deny-all policy. Explicit deny takes precedence over allow policy.

The API confirming a label does not mean that Cilium has already enforced it. Test
controlled ingress and egress peers, TCP, UDP/DNS, ICMP and established flows; correlate
the results with Cilium drop verdicts. A failed external ping alone does not establish
bidirectional isolation. See the integration checklist in the repository README.

\section{Evidence retention and retries}

The handler does not collect memory or export evidence. The default retention target
of three pods is soft: only older, managed pods with a valid quarantine timestamp and
an operator annotation asserting that evidence has been exported can be deleted.
Unknown, legacy and unexported evidence is preserved. Monitor resource consumption.

After collecting and verifying evidence, an operator can mark an individual pod:
\begin{lstlisting}[language=bash]
kubectl -n default annotate pod <quarantined-pod> \
  security-lab/evidence-exported=true
\end{lstlisting}

Deletion uses UID and resourceVersion preconditions and a grace period. Ordering is
by incident time, not pod creation time. Duplicate events reattempt cleanup without
replacing the original quarantine timestamp. There is no periodic cleanup controller.
A cleanup failure does not undo a successful quarantine: the response is still
successful, marks cleanup as pending, and the next accepted event retries it.

Falcosidekick sends each alert once and does not retry failed deliveries, so the
handler absorbs transient failures itself. Kubernetes API conflicts, throttling,
server errors and timeouts are retried up to three times with backoff, and each
attempt re-reads the pod before patching it. A target pod that no longer exists is
reported as gone. One process serializes response operations using a lock; concurrent
events wait up to a configurable time (ten seconds by default) while other threads
keep serving probes. An alert is lost only if that wait or the retries are exhausted,
in which case the handler logs an error; there is no durable queue or guaranteed
replay. Recreate prevents intended rollout overlap but causes an update interruption.
Distributed scaling requires additional design.
The target has no liveness probe that intentionally restarts isolated evidence, but
OOM and node failures can still destroy volatile state.

\section{Validation and safe cleanup}

The following checks were completed for the current implementation:
\begin{itemize}
\item Twenty-two handler regression tests pass in a hardened Python 3.11 container,
running non-root with a read-only filesystem and no network access, on every push.
\item Both container images built successfully; the target form and a loopback ping
request passed a basic smoke test.
\item Helm chart 9.1.0 rendered successfully; the webhook header and policy-selector
labels were checked in the generated resources.
\item Standard Kubernetes manifests passed server-side dry-run on Kubernetes
v1.35.6+orb1, without creating deployment resources.
\end{itemize}
These checks do not validate the Cilium custom resource or live packet enforcement.
End-to-end testing must additionally verify real Falco event identity, unauthorized
sender rejection, Service detachment, replacement readiness, bidirectional network
deny and behaviour under retries or bursts. Measure service errors and latency rather
than asserting zero downtime. The repository README contains the integration checklist.

Before cleanup, list quarantined pods, collect required evidence, and delete explicitly
identified compromised pods. Do not remove the isolation policy while compromised
pods remain running. Detached pods may outlive their Deployment; deleting the
Deployment alone is not a complete cleanup. Remove only project-owned resources and
verify the selected Kubernetes context before executing destructive commands.

\section{Results and future work}
The project connects runtime detection to an automated, conservative response.
A compromised target is removed from its Service, replaced by its ReplicaSet and
labelled for network quarantine, while it is kept running for investigation. The
handler accepts only authenticated events for an explicit set of rules, namespace
and workload; it detects conflicting writes, retries transient API failures and
never deletes evidence that an operator has not marked as exported.

The main limitations are delivery and trust. Alerts are not persisted, so a handler
outage loses them; the webhook token travels over plain HTTP inside the cluster;
and administrators, node access and label editors remain trusted. The handler does
not acquire memory or export evidence, and containment must still be measured on a
real Cilium cluster rather than inferred from accepted label changes.

Natural extensions follow from these limits:
\begin{itemize}
\item durable delivery, for example by routing Falcosidekick to a message queue
such as NATS or Kafka that the handler consumes and acknowledges;
\item transport authentication with TLS or mutual TLS between Falcosidekick and
the handler;
\item evidence acquisition before or right after isolation, for example through
the kubelet container checkpoint API, with the result recorded on the pod;
\item an automated end-to-end test on a disposable cluster that exercises each
detection scenario and verifies bidirectional network deny.
\end{itemize}
